\chapter{相关工作}\label{chap:related-work}

\section{模仿学习}

模仿学习旨在利用专家示范数据学习策略，使机器人在相似状态下产生接近专家的动作。经典行为克隆将专家轨迹分解为状态--动作样本，并通过监督学习拟合策略函数\citep{pomerleau1989alvinn,argall2009survey}。该方法实现直接、训练稳定，在示教覆盖充分、分布偏移较小的任务中效果良好。然而，机器人执行过程本质上是闭环的：策略输出会改变后续观测分布，一旦进入训练覆盖稀疏的状态区域，误差便会随时间累积并导致任务失败。Ross等人从序列预测角度分析了这一问题，指出朴素行为克隆会因训练分布与学习策略诱导分布之间的不一致而产生随时间累积的误差，并提出通过交互式数据聚合缓解分布偏移\citep{ross2010efficient,ross2011reduction}。本文采用离线示教数据，不引入在线专家查询，因此历史窗口、动作块输出和以目标为中心的表示是降低分布偏移的主要结构性机制。

动态目标抓取对模仿学习提出三个相互关联的挑战：闭环分布偏移、时间依赖性建模和动作多模态表达。针对分布偏移、时间依赖和多模态动作问题，已有研究通常从数据和模型两方面改进行为克隆。一方面，可以通过扰动采样、交互式数据聚合或合成专家数据扩展训练分布；另一方面，可以引入历史观测、目标条件和多步动作预测，使策略不再依赖单帧状态做出短视决策。确定性回归往往在一对多示教中输出条件均值，能量模型和混合密度模型等方法则常用于表达多峰动作分布\citep{bishop1994mixture,florence2022implicit}。动态目标抓取尤其需要时间信息和多模态表达，因为目标速度、相对相位（任务阶段）和夹爪闭合时机都无法由静态位置完全决定。本文比较的BC、BC-RNN、ACT和DP均属于监督模仿学习范式，但分别代表前馈回归、循环序列建模、CVAE动作分块和扩散去噪四种不同设计。

循环行为克隆是机器人离线模仿学习中的常用强基线。robomimic的基准研究表明，在多阶段操作、部分可观测和示教质量参差不齐的场景中，带循环状态的BC-RNN通常优于单帧前馈BC\citep{mandlekar2021matters}。本文的BC-RNN采用GRU作为时间状态编码器；GRU通过更新门和重置门调节历史信息流，在保持长时依赖建模能力的同时比LSTM结构更简洁\citep{cho2014learning,chung2014empirical}。这与本文在闭环控制频率下重复生成动作块的需求相符。

尽管上述方法各具优势，现有工作通常在各自的任务设置和评估条件下独立验证，缺乏在统一感知接口、统一控制器和统一噪声条件下的系统性闭环比较——这正是本文选择BC、BC-RNN、ACT和Diffusion Policy四种代表性架构进行对比研究的动机。

\section{分层模仿学习}

长时序机器人操作通常难以由单一扁平策略直接覆盖，因为策略既要理解任务意图和中间目标，又要输出稳定的连续控制。分层模仿学习由此将问题拆分为高层规划和低层执行两个部分：高层模块预测潜在计划、中间子目标或短时轨迹，低层控制器再将这些计划转化为具体动作。Play-LMP从无标签遥操作play数据中学习潜在计划空间，并在测试时复用这些计划完成用户指定的视觉操作任务\citep{lynch2020learning}；IRIS将离线示教控制问题分解为高层目标选择和低层目标条件控制器，使系统能够从大规模、质量不一的示教数据中组合更有效的执行片段\citep{mandlekar2020iris}；MimicPlay进一步利用人类play视频学习高层潜在计划，并用少量机器人遥操作数据训练低层视觉运动控制器，从而提升长时序真实机器人操作的样本效率和鲁棒性\citep{wang2023mimicplay}。然而，上述三种方法均需从数据中同时学习高层规划和低层控制，低层策略的仿真到真实迁移差距和数据覆盖问题因此被引入整个系统。这一观察促使本文将低层执行替换为模型预测控制，只保留高层规划为学习目标。

本文采用的控制器结构与上述分层范式相同，都是先由高层模块表达任务意图，再由执行层转化为物理动作；区别在于本文没有再学习一个低层运动控制器，而是让模仿策略输出以目标为中心的任务空间参考动作块，并交由MPC控制器完成约束跟踪。这种“学习规划器与MPC控制器”的设计借鉴了分层模仿学习中规划与执行解耦的思想，同时保留了MPC在关节限制、轨迹平滑和抓取后搬运等底层约束处理上的可解释性。

\section{动作分块与生成式策略}

机器人操作任务中常存在一对多映射：同一观测附近可能存在多个合理动作，例如继续跟踪、提前动态追踪或准备夹爪闭合。简单回归模型在这种情况下容易预测多个模式的平均值，造成动作保守或相位错误。动作分块方法通过一次预测未来的动作序列，将短时规划结构显式纳入策略输出\citep{zhao2023learning}，在闭环执行中有助于降低逐步误差累积。ACT由Zhao等人在低成本双臂精细操作中提出，将策略建模为条件动作序列生成，并通过重叠动作块的时间集成改善闭环平滑性\citep{zhao2023learning}。ACT类方法进一步使用Transformer建模观测上下文与动作序列内部依赖；Transformer的自注意力机制可以在无循环结构下并行建模长距离依赖\citep{vaswani2017attention}，使策略可以在固定窗口内表达较长时间的执行意图。

扩散模型为动作生成提供了另一种思路。其核心是先向数据逐步加入噪声，再训练神经网络学习反向去噪过程\citep{ho2020denoising,song2020denoising}。Diffusion Policy将这一思想用于机器人控制，把未来动作块作为条件生成样本，并通过重叠动作块的时间集成获得平滑闭环输出\citep{chi2023diffusion}。相较于确定性回归，DP原则上可以表达多峰动作分布；相较于单步策略，动作块生成更适合需要连续对齐和时机选择的操作任务。然而，生成式建模在推理效率上与确定性策略存在结构性差异：ACT和BC等确定性方法通过单次前向传播即可输出完整动作块，而DP需要迭代执行去噪采样（例如10步DDIM），每次采样均需调用去噪网络。在闭环控制中，这一成本可以被动作块执行窗口分摊：每次重规划生成的动作块覆盖未来16步，大部分控制周期仅需从时间集成缓冲区读取单步动作。尽管如此，重规划时刻的时延仍显著高于确定性策略。本文采用低维以目标为中心的状态而非图像输入，将DP与确定性动作块策略置于统一接口下比较，从而在可比的推理预算下考察不同生成机制对闭环表现的影响。

除生成机制外，动作表示方式同样显著影响模仿学习性能与稳定性。面向真实机器人泛化的Universal Manipulation Interface比较并采用了相对轨迹动作表示，以降低不同初始位姿和工作空间位置带来的数值变化\citep{chi2024universal}。Transporter Networks等目标或关键点中心方法也表明，将操作表示绑定到物体局部结构有助于获得空间泛化能力\citep{zeng2021transporter}。本文将历史观测、未来条件和动作标签统一转换到目标中心坐标系，目的也是让策略主要学习相对对齐和相位选择，而不是记忆世界坐标中的绝对轨迹。

\section{模型预测控制}

本文赋予MPC双重角色：一方面，FSM-MPC专家利用MPC生成高质量示教数据；另一方面，所有学习策略共享同一MPC控制器作为底层执行器，从而将比较焦点集中在高层参考生成策略的差异上。模型预测控制是一类滚动时域优化方法。在每个控制时刻，MPC根据系统模型预测未来状态，并求解满足约束的控制序列\citep{rawlings2017model}。在机械臂操作中，MPC能够显式考虑末端位置、姿态误差、动作平滑性、关节范围和执行速度等因素，因此常作为稳定跟踪和约束执行的核心模块。对于动态目标抓取，MPC还可以结合几何预测构造动态追踪或搬运参考轨迹。

标准MPC通常采用确定性等价假设，即认为当前状态估计是准确的，并将估计值直接代入优化问题。然而，机械臂操作中的状态估计误差不可避免：目标位姿感知可能包含高频抖动、低频漂移或偶发离群点，而这些误差会通过参考轨迹构造和优化目标函数传递到控制输出中。当目标位姿的相邻帧估计因噪声而产生虚假跳变时，MPC求解的末端参考也会随之抖动，导致夹爪闭合时机判断失准或执行轨迹不平滑。针对这一问题，现有研究提出了多种补偿方案：在控制侧，管状鲁棒MPC利用估计误差的显式界收紧约束，确保真实状态位于可行域内；在估计侧，移动时域估计和卡尔曼滤波等方法通过对历史观测序列建模来抑制高频噪声\citep{kayalibay2023filteraware}。此外，MPC的性能依赖于精确的系统动力学模型，当目标运动发生突变（如抓取过程中的阶段切换）时，其滚动优化框架无法提前做出离散决策，这正是FSM-MPC专家需要有限状态机补充阶段逻辑的原因。

本文利用MPC的双重角色，并借此考察状态估计对基于模型的方法性能的影响。在真实观测下，FSM-MPC专家作为专家策略生成高质量示教，其性能代表任务在无感知误差时的上限。在学习策略运行时，共享MPC控制器跟踪策略输出的任务空间参考，此时控制器接收的参考已经经过策略的隐式平滑，因此即使原始观测含噪声，参考序列也可能比直接使用噪声观测构造的参考更稳定。在含噪声观测实验中，本文进一步为FSM-MPC专家前接一个简单卡尔曼滤波器，以测试显式状态估计能否恢复基于模型的方法在噪声下的性能。由此，实验能够在统一控制器下区分三种情况：学习策略通过隐式时间平滑抑制噪声、未滤波的基于模型的方法将观测噪声直接传递到控制输出，以及显式状态估计帮助基于模型的方法恢复性能。

\section{动态目标抓取}

动态目标抓取要求机器人在目标运动过程中完成相对位姿对齐和接触决策，其难度高于静态抓取。抓取综述通常将解析抓取、基于采样或回归的抓取生成、强化学习以及闭环视觉伺服等路线加以区分\citep{zhang2022robotic,newbury2023deep}。传统方法通常依赖几何预测、轨迹优化、视觉伺服或显式状态机，在目标状态准确且运动模式简单时较为可靠。例如，面向移动物体的动态抓取工作利用视觉跟踪和在线轨迹修正完成动态追踪抓取\citep{marturi2019dynamic}；深度视觉伺服方法进一步将深度学习与三维视觉伺服结合用于机器人操作与抓取\citep{alshanoon2022robotic}；近年的多臂协同动态抓取研究也说明物流和工业场景中对运动目标鲁棒动态追踪的需求仍在增长\citep{he2026dynamicgrasping}。学习方法则尝试从示教数据中获得目标运动、抓取时机和控制策略之间的关系，具有适应复杂模式的潜力，但需要足够覆盖任务变化和观测扰动的数据。

本文关注平面平移与绕固定轴旋转共同作用下的动态立方体抓取。该任务同时包含连续运动跟踪、短时未来预测、阶段切换和抓取后搬运，适合作为对比基于模型的方法与模仿学习策略的实验平台。现有动态抓取研究中，解析方法在状态准确时可靠，但在噪声下退化明显。学习方法虽能适应复杂模式，但多数工作仅评估单一策略结构，缺乏对不同模仿学习架构在噪声鲁棒性上的系统比较。与直接端到端输出关节命令不同，本文采用以目标为中心的任务空间参考作为学习输出，并由MPC控制器负责低层控制。相较于§2.2中讨论的分层模仿学习方法，本文将低层执行替换为模型预测控制，形成带模型执行层的分层模仿控制框架；该设计兼顾了学习策略对动态模式的拟合能力和基于模型的方法对执行约束的可解释性，同时使不同学习策略的闭环表现可以在排除低层控制器差异的情况下进行比较\citep{lynch2020learning,mandlekar2020iris,wang2023mimicplay}。
